% =============================================================================
% 8. PRACTICAL DISCUSSION AND LIMITATIONS
% =============================================================================
\section{Discussion and Practical Limitations}
\label{sec:discussion}

\subsection{Operational Edge Deployment Considerations}
BioDINO-Edge is engineered for operational deployment on resource-constrained agricultural edge platforms, including handheld scouting terminals, tractor-mounted spray booms, and unmanned aerial vehicles (UAVs). By eliminating iterative backpropagation, the framework avoids GPU memory spikes and gradient computation graphs, executing feed-forward inference and closed-form linear algebra directly on ARM CPUs. Under prolonged continuous load, execution throughput settles at 4.26\,FPS (Section~\ref{sec:edge_hardware}); therefore, field deployment is naturally structured around shutter-triggered, on-demand inspections rather than continuous video streams. Furthermore, the modular two-tier architecture enables rapid over-the-air deployment of the lightweight base application ($\approx$22.8\,MB), allowing companion multimodal modules to be downloaded on demand when cellular bandwidth permits. In practical field scouting, agronomists can also dynamically register emerging regional pathogens through the mobile interface using 1--5 symptomatic foliar photos and an optional healthy reference leaf, with closed-form parameter adaptation executing directly on the device CPU.

\subsection{Practical Limitations and Scope}
\textbf{Canopy clutter and leaf segmentation.} BioDINO-Edge processes foliar images at the individual leaf frame level. In unsegmented field imagery with dense, overlapping canopies, multiple leaves from neighboring plants can enter the camera field of view, lowering gating confidence. While tight lesion cropping degrades host context, processing unsegmented canopy clutter introduces irrelevant background vectors. Future work could incorporate leaf segmentation to isolate individual leaves before closed-form classification.

\textbf{Statistical pseudo-replication in support expansion.} Multi-view support expansion generates correlated pseudo-replicates. Because the analytical Ledoit--Wolf shrinkage intensity assumes independent and identically distributed (i.i.d.)\ observations, the nominal $N-C=429$ degrees of freedom overstate the independent information in the support set (effective sample size: 55--440). Although the sensitivity study in Section~\ref{sec:aug_sensitivity} shows that downstream accuracy remains stable across view counts and augmentation modes, a cluster-aware shrinkage estimator that explicitly accounts for within-image dependence remains an important theoretical direction.

\textbf{Attention artifacts in frozen backbones.} Frozen vision transformers can produce high-norm artifact tokens in background regions, and these tokens distort attention maps~\cite{darcet2024vision}. We see similar artifacts on non-symptomatic leaf structures such as veins and petioles. Because Fisher-LDA operates on pooled representations downstream, it does not alter internal backbone attention; the framework targets whole-leaf diagnosis rather than pixel-level lesion localization.

\textbf{Optical degradation and focus blur.} Severe defocus or motion blur degrades fine lesion textures, reducing diagnostic accuracy. The Android client screens frames before diagnosis with two lightweight guards: the maximum cosine similarity between the DINOv2 embedding and the class prototypes must exceed 0.12 for the frame to count as foliage, and the mean luminance of a 64$\times$64 center crop must exceed 22/255 for it to count as illuminated. If either check fails, the application suppresses the diagnosis to prevent false predictions on uninformative frames. An additional discrete-Laplacian variance check (threshold 45) detects motion blur, prompting the user to hold the camera steady when frames are blurred rather than issuing an uncertain diagnosis. While the uncertainty gate routes ambiguous clear foliage to Tier-2, handling optical degradation relies on these front-end guards.

\textbf{Scope and threats to validity.} The physical edge profiling in this work was conducted on a single physical device (Samsung Galaxy S24~FE) with CPU-only INT8 ONNX Runtime; absolute latency and memory will vary across mobile SoCs, though the underlying algorithmic complexity remains hardware-independent. In few-shot evaluation, BioDINO-Edge is designed specifically for streaming, sequential edge inference rather than batch transduction over pre-collected query pools. Finally, while PlantDoc evaluates robustness on unconstrained field foliage with target-domain support, fully unsupervised cross-domain transfer (e.g., PlantVillage$\rightarrow$PlantDoc without target support), quantitative image-quality stress tests, and multi-device profiling remain future research directions.

% =============================================================================
